\origpage{47}{GRAPH MINORS. XIV}
Now by (6.2), $bd(\bar{r}) = U(C)$ for some circuit $C$ of $H$; let $Q$ be the path of $C$ between $u$ and $v$ with $U(Q) \subseteq \Delta$. Then $P \cup Q$ is a circuit. Let $C'_1 = C_1, \dots, C_{\beta+1} = C'_{\beta+1}$ except that $C'_h$ is obtained from $C_h$ by replacing $P$ by $Q$. Let $P'_1 = P_1, \dots, P'_\beta = P_\beta$ except that if $P$ is a subpath of some $P_j$ we replace $P$ by $Q$. Since none of $U(C_1), \dots, U(C_{\beta+1})$ intersects $\Delta$ except $U(C_h)$, it follows that $C'_1, \dots, C'_{\beta+1}$ satisfy (1). Since none of $U(P_1), \dots, U(P_\beta)$ meet $\Delta - U(P)$ (because each intersects $C_h$ in a path) it follows that $P'_1, \dots, P'_\beta$ satisfy (2). But also (3) is satisfied, as we can see as follows. We only need consider the intersections of $P'_1, \dots, P'_\beta$ with $C'_h$. Thus, let $1 \le j \le \beta$. If $P$ is a subpath of $P_j$ then $P'_j \cap C'_h$ is obtained from $P_j \cap C_h$ by replacing $P$ by $Q$, and so we may assume that it is not. If $P \cap P_j$ is null then $P'_j \cap C'_h = P_j \cap C_h$. We may assume then that $P \cap P_j$ is a path with one end $v_1$ or $v_2$ and the other internal to $P$. (Not both ends are internal to $P$, because only one edge of $P_j$ in $\Delta_h - bd(\Delta_h)$ is incident with a vertex of $C_h$ and $U(P_j) \cap (\Delta - U(P)) = \emptyset$.) But then $P'_j \cap C'_h$ is obtained from the path $P_j \cap C_h$ by deleting a terminal subpath, and so $P'_j \cap C'_h$ is a path. Thus (3) is satisfied. But this contradicts the maximality of $\Delta_h$, and so proves (4).

We see that $\Lambda = \Delta_h - bd(\Delta_h)$ is a $(3\beta + 3)$-zone around $z$, because $\Lambda \subseteq \Lambda_1$. Let $\mathscr{T}'$ be obtained by clearing $\Lambda$, where $\mathscr{T}'$ has order $\theta - 12\beta - 14$. By (4) and (6.3), $H \cap (\Sigma - \Lambda)$ is rigid. For $1 \le j \le \beta$, let $v_j$ be the first vertex of $P_j$ (starting from $bd(\Lambda_2)$) which is in $V(C_h)$, and let $X = \{v_1, \dots, v_\beta\}$. Then there are $\beta$ paths of $H$ from $V(C_1)$ to $X$, mutually disjoint, and each meeting $C_h$ just its last vertex (namely, the subpath of $P_j$ from $V(C_1)$ to $v_j$, for $1 \le j \le \beta$). Moreover, $\Lambda$ is a $(3\beta + 3)$-zone around $z$, and $\Lambda_1$ is a $(\theta - \beta - 4(3\beta + 3) - 2)$-zone around $z$, since $\theta - \beta - 4(3\beta + 3) - 2 \ge 3\beta + 3$. Thus, by (7.8) (applied to $H \cap bd(\Lambda_1)$ and $C_h$, and with $\lambda = 3\beta + 3$), it follows that $X$ is free with respect to $\mathscr{T}'$. Hence $(\Lambda, C, X)$ satisfies the theorem. $\square$

If $\alpha$ is an arrangement in a surface $\Sigma$ of a segregation $\mathscr{S}$ of $G$, and $\Sigma' \subseteq \Sigma$ is a surface with $bd(\Sigma')$ $\alpha$-normal, then for each $(A, \Omega) \in \mathscr{S}$, either $\alpha(A, \Omega) \subseteq \Sigma'$ or $\alpha(A, \Omega) \cap \Sigma' \subseteq bd(\Sigma') \cap bd(\alpha(A, \Omega))$. We define $\alpha^{-1}(\Sigma')$ to be the subgraph of $G$ formed by
\begin{equation*}
\bigcup (A : (A, \Omega) \in \mathscr{S} \text{ and } \alpha(A, \Omega) \subseteq \Sigma')
\end{equation*}
together with all $v \in V(\mathscr{S})$ such that $\alpha(v) \in \Sigma'$. We apply (9.5) to prove the following.

\stmt
(9.6) \textit{Let $H$ be a 2-cell, rigid drawing in $\Sigma$, let $z \in A(H)$, and let $\mathscr{T}$ be a respectful tangle in $H$ of order $\theta$, with metric $d$. Let $\pi \ge 0$ with $\pi \ne 1, 3$, let $\beta \ge 2\pi + 8$ be even, and let $\theta \ge 16\beta + 17$. Let $H$ be a subgraph of $G$, let $J$ be the union of some bridges of $H$ in $G$ such that $d(z, v) < \beta$ for each}
